\documentclass{article}
\usepackage[utf8]{inputenc}
\usepackage{array}
\usepackage[margin=2cm]{geometry}

\newcolumntype{C}[1]{>{\centering\arraybackslash}m{#1}}
\newcolumntype{L}[1]{>{\raggedright\arraybackslash}m{#1}}

\begin{document}

\begin{center}
{\Large \textbf{REQUERIMIENTOS FUNCIONALES}}\\[4pt]
Sistema de Gestión de Inventario Multi-canal --- Versión 1.0
\end{center}

\section*{RF-01: Registrar producto}
\begin{center}
\renewcommand{\arraystretch}{1.4}
\begin{tabular}{|L{4cm}|L{12cm}|}
\hline
\textbf{ID} & RF-01 \\
\hline
\textbf{Nombre} & Registrar producto \\
\hline
\textbf{Descripción} & El sistema debe permitir registrar un producto con nombre, categoría, precio unitario, cantidad inicial y stock mínimo. \\
\hline
\textbf{Actor} & Administrador \\
\hline
\textbf{Entradas} & Nombre, Categoría, Precio unitario, Cantidad inicial, Stock mínimo \\
\hline
\textbf{Salidas} & Producto registrado en la base de datos, Movimiento de entrada inicial generado \\
\hline
\textbf{Precondiciones} & El administrador debe estar autenticado, Debe existir al menos una categoría registrada \\
\hline
\textbf{Postcondiciones} & El producto queda disponible para registrar entradas y salidas, El producto aparece en la lista de productos \\
\hline
\textbf{Prioridad} & Alta \\
\hline
\textbf{CU relacionado} & CU-002 \\
\hline
\end{tabular}
\end{center}

\section*{RF-02: Registrar entrada de stock}
\begin{center}
\renewcommand{\arraystretch}{1.4}
\begin{tabular}{|L{4cm}|L{12cm}|}
\hline
\textbf{ID} & RF-02 \\
\hline
\textbf{Nombre} & Registrar entrada de stock \\
\hline
\textbf{Descripción} & El sistema debe permitir registrar una entrada de stock (compra a proveedor) indicando producto, cantidad y nombre del proveedor como dato descriptivo. \\
\hline
\textbf{Actor} & Administrador \\
\hline
\textbf{Entradas} & Producto, Cantidad, Nombre del proveedor \\
\hline
\textbf{Salidas} & Stock del producto incrementado, Movimiento de entrada registrado con fecha y hora \\
\hline
\textbf{Precondiciones} & El administrador debe estar autenticado, El producto debe existir en el sistema \\
\hline
\textbf{Postcondiciones} & El stock del producto queda actualizado, El movimiento queda registrado en el historial \\
\hline
\textbf{Prioridad} & Alta \\
\hline
\textbf{CU relacionado} & CU-003 \\
\hline
\end{tabular}
\end{center}

\section*{RF-03: Registrar salida de stock}
\begin{center}
\renewcommand{\arraystretch}{1.4}
\begin{tabular}{|L{4cm}|L{12cm}|}
\hline
\textbf{ID} & RF-03 \\
\hline
\textbf{Nombre} & Registrar salida de stock \\
\hline
\textbf{Descripción} & El sistema debe permitir registrar una salida de stock (venta) indicando producto, cantidad y canal de venta (físico, WhatsApp, redes sociales o tienda virtual). \\
\hline
\textbf{Actor} & Administrador / Vendedor \\
\hline
\textbf{Entradas} & Producto, Cantidad, Canal de venta \\
\hline
\textbf{Salidas} & Stock del producto decrementado, Movimiento de salida registrado con fecha y hora \\
\hline
\textbf{Precondiciones} & El usuario debe estar autenticado, El producto debe existir y tener stock disponible \\
\hline
\textbf{Postcondiciones} & El stock del producto queda actualizado, El movimiento queda registrado en el historial \\
\hline
\textbf{Prioridad} & Alta \\
\hline
\textbf{CU relacionado} & CU-004 \\
\hline
\end{tabular}
\end{center}

\section*{RF-04: Actualizar stock automáticamente}
\begin{center}
\renewcommand{\arraystretch}{1.4}
\begin{tabular}{|L{4cm}|L{12cm}|}
\hline
\textbf{ID} & RF-04 \\
\hline
\textbf{Nombre} & Actualizar stock automáticamente \\
\hline
\textbf{Descripción} & El sistema debe actualizar automáticamente la cantidad disponible de un producto al registrar una entrada o salida. \\
\hline
\textbf{Actor} & Sistema \\
\hline
\textbf{Entradas} & Movimiento de entrada o salida confirmado \\
\hline
\textbf{Salidas} & Cantidad disponible del producto actualizada \\
\hline
\textbf{Precondiciones} & Debe haberse registrado una entrada (RF-02) o una salida (RF-03) \\
\hline
\textbf{Postcondiciones} & El stock refleja el valor correcto tras el movimiento, Se recalcula el estado de stock (disponible/bajo/agotado) \\
\hline
\textbf{Prioridad} & Alta \\
\hline
\textbf{CU relacionado} & CU-002, CU-003, CU-004 \\
\hline
\end{tabular}
\end{center}

\section*{RF-05: Mostrar alerta visual de stock bajo}
\begin{center}
\renewcommand{\arraystretch}{1.4}
\begin{tabular}{|L{4cm}|L{12cm}|}
\hline
\textbf{ID} & RF-05 \\
\hline
\textbf{Nombre} & Mostrar alerta visual de stock bajo \\
\hline
\textbf{Descripción} & El sistema debe mostrar una alerta visual en el panel principal cuando el stock de un producto sea igual o inferior al mínimo configurado. La notificación por email o mensajería queda fuera del alcance de la primera versión. \\
\hline
\textbf{Actor} & Sistema \\
\hline
\textbf{Entradas} & Stock actual del producto, Stock mínimo configurado \\
\hline
\textbf{Salidas} & Alerta visual mostrada en el panel principal \\
\hline
\textbf{Precondiciones} & El producto debe tener configurado un stock mínimo, Debe existir al menos un movimiento registrado \\
\hline
\textbf{Postcondiciones} & El producto aparece marcado visualmente en el panel, La alerta queda registrada con fecha y hora \\
\hline
\textbf{Prioridad} & Alta \\
\hline
\textbf{CU relacionado} & CU-005 \\
\hline
\end{tabular}
\end{center}

\section*{RF-06: Consultar lista de productos}
\begin{center}
\renewcommand{\arraystretch}{1.4}
\begin{tabular}{|L{4cm}|L{12cm}|}
\hline
\textbf{ID} & RF-06 \\
\hline
\textbf{Nombre} & Consultar lista de productos \\
\hline
\textbf{Descripción} & El sistema debe permitir consultar la lista completa de productos con su stock actual. \\
\hline
\textbf{Actor} & Administrador / Vendedor \\
\hline
\textbf{Entradas} & Solicitud de consulta de productos \\
\hline
\textbf{Salidas} & Lista de productos con nombre, categoría, precio, stock y estado \\
\hline
\textbf{Precondiciones} & El usuario debe estar autenticado \\
\hline
\textbf{Postcondiciones} & Se muestra la información sin modificar el sistema \\
\hline
\textbf{Prioridad} & Alta \\
\hline
\textbf{CU relacionado} & CU-006 \\
\hline
\end{tabular}
\end{center}

\section*{RF-07: Filtrar productos}
\begin{center}
\renewcommand{\arraystretch}{1.4}
\begin{tabular}{|L{4cm}|L{12cm}|}
\hline
\textbf{ID} & RF-07 \\
\hline
\textbf{Nombre} & Filtrar productos \\
\hline
\textbf{Descripción} & El sistema debe permitir filtrar productos por categoría, nombre y estado de stock, donde: Disponible = stock $>$ stock mínimo; Bajo = 0 $<$ stock $\leq$ stock mínimo; Agotado = stock = 0. \\
\hline
\textbf{Actor} & Administrador / Vendedor \\
\hline
\textbf{Entradas} & Categoría, Nombre, Estado de stock \\
\hline
\textbf{Salidas} & Lista filtrada de productos \\
\hline
\textbf{Precondiciones} & El usuario debe estar autenticado, Debe existir al menos un producto registrado \\
\hline
\textbf{Postcondiciones} & Se muestra la lista filtrada sin modificar el sistema \\
\hline
\textbf{Prioridad} & Media \\
\hline
\textbf{CU relacionado} & CU-007 \\
\hline
\end{tabular}
\end{center}

\section*{RF-08: Consultar historial de movimientos por producto}
\begin{center}
\renewcommand{\arraystretch}{1.4}
\begin{tabular}{|L{4cm}|L{12cm}|}
\hline
\textbf{ID} & RF-08 \\
\hline
\textbf{Nombre} & Consultar historial de movimientos por producto \\
\hline
\textbf{Descripción} & El sistema debe permitir consultar el historial de movimientos de un producto específico (entradas y salidas), ordenado cronológicamente. \\
\hline
\textbf{Actor} & Administrador \\
\hline
\textbf{Entradas} & Producto seleccionado \\
\hline
\textbf{Salidas} & Listado cronológico de movimientos: fecha, tipo, cantidad, proveedor o canal \\
\hline
\textbf{Precondiciones} & El usuario debe estar autenticado, El producto debe existir y tener al menos un movimiento registrado \\
\hline
\textbf{Postcondiciones} & Se muestra el historial sin modificar el sistema \\
\hline
\textbf{Prioridad} & Media \\
\hline
\textbf{CU relacionado} & CU-008 \\
\hline
\end{tabular}
\end{center}

\section*{RF-09: Consultar estadísticas básicas}
\begin{center}
\renewcommand{\arraystretch}{1.4}
\begin{tabular}{|L{4cm}|L{12cm}|}
\hline
\textbf{ID} & RF-09 \\
\hline
\textbf{Nombre} & Consultar estadísticas básicas \\
\hline
\textbf{Descripción} & El sistema debe mostrar estadísticas agregadas del negocio: productos más vendidos en un período y resumen global de movimientos del período (no por producto). \\
\hline
\textbf{Actor} & Administrador \\
\hline
\textbf{Entradas} & Período de consulta \\
\hline
\textbf{Salidas} & Ranking de productos más vendidos, Resumen de movimientos del período \\
\hline
\textbf{Precondiciones} & El usuario debe estar autenticado, Deben existir movimientos registrados en el período \\
\hline
\textbf{Postcondiciones} & Se muestran las estadísticas sin modificar el sistema \\
\hline
\textbf{Prioridad} & Media \\
\hline
\textbf{CU relacionado} & CU-009 \\
\hline
\end{tabular}
\end{center}

\section*{RF-10: Gestionar usuarios}
\begin{center}
\renewcommand{\arraystretch}{1.4}
\begin{tabular}{|L{4cm}|L{12cm}|}
\hline
\textbf{ID} & RF-10 \\
\hline
\textbf{Nombre} & Gestionar usuarios \\
\hline
\textbf{Descripción} & El sistema debe permitir gestionar usuarios con diferentes roles (administrador y vendedor). \\
\hline
\textbf{Actor} & Administrador \\
\hline
\textbf{Entradas} & Nombre de usuario, Contraseña, Rol (Administrador / Vendedor), Estado (activo / inactivo) \\
\hline
\textbf{Salidas} & Usuario registrado o actualizado en la base de datos \\
\hline
\textbf{Precondiciones} & El administrador debe estar autenticado \\
\hline
\textbf{Postcondiciones} & El usuario queda habilitado para iniciar sesión según su rol \\
\hline
\textbf{Prioridad} & Alta \\
\hline
\textbf{CU relacionado} & CU-001, CU-010 \\
\hline
\end{tabular}
\end{center}

\section*{RF-11: Restringir acceso a configuración}
\begin{center}
\renewcommand{\arraystretch}{1.4}
\begin{tabular}{|L{4cm}|L{12cm}|}
\hline
\textbf{ID} & RF-11 \\
\hline
\textbf{Nombre} & Restringir acceso a configuración \\
\hline
\textbf{Descripción} & El sistema debe restringir las funciones de gestión de usuarios, categorías y configuración del sistema al rol administrador. El rol vendedor podrá registrar salidas (RF-03) y consultar productos (RF-06, RF-07), pero no gestionar entradas ni configuración. \\
\hline
\textbf{Actor} & Sistema \\
\hline
\textbf{Entradas} & Solicitud de acceso a función restringida, Rol del usuario \\
\hline
\textbf{Salidas} & Acceso permitido o denegado, Mensaje ``No tiene permisos para esta función'' \\
\hline
\textbf{Precondiciones} & El usuario debe estar autenticado \\
\hline
\textbf{Postcondiciones} & El acceso queda controlado según el rol, El intento no autorizado queda registrado en log \\
\hline
\textbf{Prioridad} & Alta \\
\hline
\textbf{CU relacionado} & CU-001, CU-010, CU-011 \\
\hline
\end{tabular}
\end{center}

\section*{RF-12: Gestionar categorías de productos}
\begin{center}
\renewcommand{\arraystretch}{1.4}
\begin{tabular}{|L{4cm}|L{12cm}|}
\hline
\textbf{ID} & RF-12 \\
\hline
\textbf{Nombre} & Gestionar categorías de productos \\
\hline
\textbf{Descripción} & El sistema debe permitir registrar y gestionar categorías de productos (crear, editar y eliminar). No se podrá eliminar una categoría con productos asociados. \\
\hline
\textbf{Actor} & Administrador \\
\hline
\textbf{Entradas} & Nombre de la categoría \\
\hline
\textbf{Salidas} & Categoría registrada, actualizada o eliminada en el sistema \\
\hline
\textbf{Precondiciones} & El administrador debe estar autenticado \\
\hline
\textbf{Postcondiciones} & La categoría queda disponible para asignar a productos \\
\hline
\textbf{Prioridad} & Media \\
\hline
\textbf{CU relacionado} & CU-012 \\
\hline
\end{tabular}
\end{center}

\end{document}
